% Procedencia literal: output/PREPARACION_ENTREGA_USB_20260928/ENTREGA_HMT_USB/03_FUENTES_Y_REPRODUCCION/ACTUALIZACIONES_20260928/FINAL01/PAQUETES/ARTICULO_XI_REV11_ES_EN_FINAL_BN_20260928/ARTICULO_XI_REV11_ES_EN_FINAL_BN/ES/source/sections/ch_tres_vueltas_ext.tex
% Recuperación íntegra; única transformación mecánica: prefijo hmtfg: en label/ref/eqref.
% No se modifican operadores, pruebas, tablas ni afirmaciones del propietario.
\subsubsection{Realización arista por arista de las tres prolongaciones nonádicas}
\label{hmtfg:subsubsec:tres-vueltas-ext-ch}

Esta subsección especifica el modelo autónomo generado por la relación de
extensión TPK que interviene en la prueba cardinal. No se infiere la existencia de
una arista por coincidencia de tipos: se construyen sus estados inicial y final,
su relación APP, su transporte de fibra y la actualización unaria del registro.
La construcción utiliza exclusivamente las dos hojas de APP, el selector TRIT
y la composición \(\mathsf{Sel}\)--\(\mathsf{Tra}\)--\(\operatorname{Upd}\) del TPK.

\paragraph{Células completas y tres reducciones balanceadas.}
En la carta \(D_9=\{1,\ldots,9\}\), mantendremos siempre la célula APP
completa
\[
 a_{p,q}:=(p,q;R^\Sigma_{p,q},q^\Sigma_{p,q},
                 R^\Pi_{p,q},q^\Pi_{p,q}),
\]
con
\[
 p+q=R^\Sigma_{p,q}+9q^\Sigma_{p,q},\qquad
 pq=R^\Pi_{p,q}+9q^\Pi_{p,q}.
\]
La hoja activa se registra en una coordenada separada; no se escinde
\(a_{p,q}\) en una supuesta célula aditiva y otra multiplicativa. Para
\(r\in D_9\), sean
\begin{equation}
 a(r):=
 \begin{cases}
  a_{1,9},&r=1,\\
  a_{1,r-1},&2\le r\le9,
 \end{cases}
 \qquad
 \tau(r):=M_{\rm ph}(r),
 \qquad
 m(r):=\tau(r)\in\{-1,0,+1\}.
 \label{hmtfg:eq:celula-completa-fase-ch}
\end{equation}
El residuo aditivo de \(a(r)\) proporciona aquí un representante APP de la
fase; esto no convierte \(\Sigma\) en hoja activa universal. El selector
dinámico activa \(\Sigma\) en \(r\in\{1,4,7\}\), activa \(\Pi\) en
\(r\in\{2,5,8\}\) y conserva ambas hojas, sin desplazar cursor, en
\(r\in\{3,6,9\}\). En los tres casos se transportan conjuntamente
\(R^\Sigma,q^\Sigma,R^\Pi,q^\Pi\). Se tiene
\[
 (m(1),\ldots,m(9))=(1,-1,0,1,-1,0,1,-1,0).
\]
Las tres parejas ordenadas por la base nonádica son
\begin{equation}
 \mathfrak p_{+1}:=(1,4),\qquad
 \mathfrak p_{-1}:=(2,5),\qquad
 \mathfrak p_{0}:=(3,6).
 \label{hmtfg:eq:tres-pares-internos-ch}
\end{equation}
Estas parejas son la sección normalizada que toma las dos primeras
ocurrencias de cada tipo TRIT bajo el desplazamiento \(r\mapsto r+3\); el
tercer bloque de fases \(7,8,9\) conserva el control de cierre de la vuelta.
Sus dos miembros pertenecen a una misma fibra de \(M_{\rm ph}\). La
reducción balanceada definida en estos antecedentes da, sin introducir un coeficiente
objetivo,
\begin{equation}
 \begin{array}{c|c|c|c}
 \varepsilon&m(\mathfrak p_\varepsilon^{(1)})
                  +m(\mathfrak p_\varepsilon^{(2)})
   &r_3\bigl(2\varepsilon\bigr)&
   \chi(\varepsilon,\varepsilon)\\ \hline
 +1&1+1=(-1)+3\cdot1&-1&+1\\
 0&0+0=0+3\cdot0&0&0\\
 -1&(-1)+(-1)=1+3\cdot(-1)&+1&-1
 \end{array}
 \label{hmtfg:eq:tres-carry-derivados-ch}
\end{equation}
Por tanto,
\begin{equation}
 \chi(\varepsilon,\varepsilon)
 =\frac{2\varepsilon-r_3(2\varepsilon)}3=\varepsilon.
 \label{hmtfg:eq:carry-no-parametrico-ch}
\end{equation}
El signo de la rama es así la salida del cociclo balanceado aplicado a dos
emisiones TRIT producidas por APP. Cada pareja proporciona una de las tres
salidas de acarreo. Por claridad, indexaremos después la pareja mediante
\(\varepsilon(\mathfrak p):=
\chi(m(\mathfrak p^{(1)}),m(\mathfrak p^{(2)}))\); el símbolo
\(\varepsilon\) no suministra a \(\operatorname{Upd}^{\rm cal}\) un acarreo libre.

La arquitectura de fase, el levantamiento balanceado y el transporte de
memoria se realizan aquí mediante veintisiete incidencias locales y sus
grafos de extensión. La sección normalizada fija representantes de las
fibras de \(M_{\rm ph}\); la cobertura que se demostrará utiliza la existencia
de esos representantes y permanece invariante al sustituirlos por otra
sección compatible.

\paragraph{Datos tipados de las nueve aristas.}
Sea
\(K_{\rm Carr}:=\langle\kappa_{\rm Carr}\rangle_{\rm Ab}\) el módulo de
acarreo. Las acciones están tipadas por
\[
 H_{\rm Carr}:\mathscr P_{\rm TPK}^{\rm op}
   \longrightarrow\operatorname{End}(K_{\rm Carr}),\qquad
 Y_{\rm Carr}:\mathscr P_{\rm TPK}
   \longrightarrow\operatorname{End}(K_{\rm Carr}).
\]
En este modelo generado se toman las acciones triviales
\(H_{\rm Carr}(e)=Y_{\rm Carr}(e)=\operatorname{id}_{K_{\rm Carr}}\): la
hoja activa y la ruta se conservan como coordenadas propias del registro y no
actúan sobre el sumando libre de acarreo. Sea además
\((\mathsf M_d,\jmath_{d',d})\) el sistema directo graduado de memoria. Para
no confundir dos representantes identificados por el colímite, el estado
conserva además el grado como coordenada y se usa el portador etiquetado
\[
 \widetilde{\mathsf M}:=
 \bigsqcup_{d\ge1}\bigl(\{d\}\times\mathsf M_d\bigr),
 \qquad
 \operatorname{gr}_{\mathsf M}(d,\mu)=d.
\]
La publicación en la memoria límite viene dada por
\[
\begin{aligned}
 \mathsf M&:=\varinjlim_d\mathsf M_d,&
 \jmath_d&:\mathsf M_d\longrightarrow\mathsf M,\\
 u&:=\jmath_d(u_d),&
 \iota_{\mathsf M}:\mathbb Z&\longrightarrow\mathsf M,
 \qquad n\longmapsto nu .
\end{aligned}
\]
Aquí \(\jmath_d\) es el mapa canónico del colímite y los \(u_d\) forman una
familia compatible. Su componente libre permite
el lector \(\deg_u(nu)=n\).
Las proyecciones de configuración, fase, memoria, acarreo, ruta y
supervivencia se tiparán sobre el portador calibrado construido
recursivamente más abajo; no se importará para ellas ningún espacio ambiental.
Para cada arista construida \(e\), la
componente calibrada
\(\operatorname{Upd}^{\rm cal}_e:=\mathsf{Mem}^{\rm cal}_e\) denotará la
tercera aplicación de la factorización; su ley se fija campo por campo más
abajo y no se presupone como restricción de una relación ambiental.

Fijemos la profundidad inicial \(k_0=1\), una configuración completa
\(q_\star\in\mathcal Q_{\rm obs}\) de fase uno y
\[
 q_{n,r}:=F_{\rm obs}^{9n+r-1}(q_\star),\qquad
 q_{n,10}:=q_{n+1,1}.
\]
Como el calendario observable es periódico, su valor \(q_{n,r}\) no
distingue por sí solo apariciones separadas por doce vueltas. El modelo
calibrado emplea, por ello, la elevación etiquetada
\[
 \widetilde q_{n,r}:=(q_{n,r},n,r),\qquad
 \widetilde q_{n,10}:=\widetilde q_{n+1,1},\qquad
 \pi_{\rm obs}(\widetilde q_{n,r})=q_{n,r}.
\]
La dinámica elevada
\(\widetilde F_{\rm obs}(\widetilde q_{n,r})
=\widetilde q_{n,r+1}\) proyecta a \(F_{\rm obs}\); la etiqueta no modifica
la fase física, sino que hace distintos sus niveles de aparición.
Se etiqueta la copia de la célula en cada nivel mediante
\[
 a_{n,r}:=(n,r;a(r)),\qquad
 (n,r)^+:=
 \begin{cases}
  (n,r+1),&1\le r<9,\\
  (n+1,1),&r=9.
 \end{cases}
\]
La etiqueta sólo evita identificar dos apariciones de una misma célula; las
dos divisiones euclídeas y sus cuatro coordenadas son las de \(a(r)\).
Para la vuelta \(n\ge0\) y la fase \(r\in D_9\), el símbolo
\[
 e_{\varepsilon,n,r}:\widetilde q_{n,r}
 \longrightarrow\widetilde q_{n,r+1},
 \qquad r^+:=1+(r\bmod9),
\]
ocupa el nivel
\(d_{n,r}:=k_n+r-1\), donde \(k_n:=k_0+9n\). Su registro local contiene
la célula completa \(a(r)\), la hoja activa determinada por \(M_{\rm ph}(r)\),
el tipo \(\tau(r)\),
la marca de pertenencia a \(\mathfrak p_\varepsilon\) y, al alcanzar el segundo miembro
de esa pareja, la identidad certificada
\eqref{hmtfg:eq:tres-carry-derivados-ch}. Definimos sus operadores únicamente a
partir de ese registro:
\begin{align}
 \mathsf C_{\mathsf M}(e_{\varepsilon,n,r})
   &=\jmath_{d_{n,r}+1,d_{n,r}},&
 \kappa_{\mathsf M}(e_{\varepsilon,n,r})
   &=\mathbf1_{\{r=9\}}u_{d_{n,r}+1},
 \label{hmtfg:eq:memoria-arista-ch}\\
       \operatorname{Carr}(e_{\varepsilon,n,r})
   &=d_{\varepsilon,r}\kappa_{\rm Carr},&
 d_{\varepsilon,r}
   &:=\mathbf1_{\{r=\mathfrak p_\varepsilon^{(2)}\}}
       \chi\!\left(m(\mathfrak p_\varepsilon^{(1)}),
                    m(\mathfrak p_\varepsilon^{(2)})\right).
 \label{hmtfg:eq:carry-arista-derivado-ch}
\end{align}
Para las identidades y para toda composición admisible se define, además,
\begin{align}
 \mathsf C_{\mathsf M}(\operatorname{id}_{\widetilde q};d)
   &:=\operatorname{id}_{\mathsf M_d},&
 \kappa_{\mathsf M}(\operatorname{id}_{\widetilde q};d)&:=0,
 \label{hmtfg:eq:memoria-identidad-ch}\\
 \mathsf C_{\mathsf M}(e_t\cdots e_1)
   &:=\mathsf C_{\mathsf M}(e_t)\circ\cdots\circ
      \mathsf C_{\mathsf M}(e_1),&
 \kappa_{\mathsf M}(e_2e_1)
   &:=\mathsf C_{\mathsf M}(e_2)
        \bigl(\kappa_{\mathsf M}(e_1)\bigr)
      +\kappa_{\mathsf M}(e_2).
 \label{hmtfg:eq:cociclo-memoria-calibrado-ch}
\end{align}
La identidad lleva el grado \(d\) porque una misma configuración observable
puede reaparecer en profundidades distintas. La última igualdad, iterada, define
\(\kappa_{\mathsf M}(e_t\cdots e_1)\) para cualquier longitud y es la ley de
cociclo del sistema directo; las inclusiones satisfacen
\(\jmath_{d+2,d+1}\jmath_{d+1,d}=\jmath_{d+2,d}\).
Además se fija
\(s(e_{\varepsilon,n,r})=\top\). Las siguientes veintisiete entradas
certifican las coordenadas que distinguen las tres prolongaciones; los campos
universales del transporte se especifican inmediatamente después. Los símbolos \(I\) y \(II\)
designan, respectivamente, la primera emisión conservada y la segunda emisión
que efectúa la reducción; un guion indica incremento de acarreo nulo, pero no
ausencia de arista.
\begin{center}
\small
\setlength{\tabcolsep}{3.2pt}
\begin{tabular}{c c cc cc cc c}
\toprule
\(r\)&\(M_{\rm ph}(r)\)&\multicolumn{2}{c}{\(\mathfrak p_{+1}\)}&
\multicolumn{2}{c}{\(\mathfrak p_0\)}&
\multicolumn{2}{c}{\(\mathfrak p_{-1}\)}&
\(\Delta\operatorname{mem}\)\\
& &marca&\(d_{+1,r}\)&marca&\(d_{0,r}\)&marca&\(d_{-1,r}\)&\\
\midrule
1&\(+1\)&I&0&--&0&--&0&0\\
2&\(-1\)&--&0&--&0&I&0&0\\
3&0&--&0&I&0&--&0&0\\
4&\(+1\)&II&\(+1\)&--&0&--&0&0\\
5&\(-1\)&--&0&--&0&II&\(-1\)&0\\
6&0&--&0&II&0&--&0&0\\
7&\(+1\)&--&0&--&0&--&0&0\\
8&\(-1\)&--&0&--&0&--&0&0\\
9&0&--&0&--&0&--&0&\(u\)\\
\bottomrule
\end{tabular}
\end{center}
Así, \(d_{\varepsilon,r}\) se calcula a partir de dos salidas del selector y
es distinto de los cocientes \(q^\Sigma,q^\Pi\) de la célula. El registro
conserva también el par marcado y el residuo balanceado; en particular, la
rama nula no se confunde con una ausencia de evento.

La ley de composición del acarreo empleada por la categoría TPK es
\begin{equation}
 \operatorname{Carr}(\operatorname{id}_{\widetilde q})=0,\qquad
 \operatorname{Carr}(e_2e_1)
 =H_{\rm Carr}(e_1)\bigl(\operatorname{Carr}(e_2)\bigr)
  +Y_{\rm Carr}(e_2)\bigl(\operatorname{Carr}(e_1)\bigr).
 \label{hmtfg:eq:ley-carry-correcta-ch}
\end{equation}
En identidades y caminos se prolongan functorialmente:
\[
\begin{aligned}
 H_{\rm Carr}(\operatorname{id})
   &=Y_{\rm Carr}(\operatorname{id})=\operatorname{id},\\
 H_{\rm Carr}(e_t\cdots e_1)
   &=H_{\rm Carr}(e_1)\circ\cdots\circ H_{\rm Carr}(e_t),\\
 Y_{\rm Carr}(e_t\cdots e_1)
   &=Y_{\rm Carr}(e_t)\circ\cdots\circ Y_{\rm Carr}(e_1).
\end{aligned}
\]
En este modelo todas estas acciones siguen siendo la identidad, pero
su tipado explicita qué endomorfismo actúa sobre cada sumando de
\eqref{hmtfg:eq:ley-carry-correcta-ch}.
En las aristas calibradas, \(H_{\rm Carr}=Y_{\rm Carr}=\operatorname{id}\),
por lo que la composición suma los incrementos ya producidos. La ecuación
\eqref{hmtfg:eq:ley-carry-correcta-ch}, y no una identificación entre orientación y
cociente, gobierna el registro acumulado.

\paragraph{Construcción independiente de las tuplas de transición.}
Un estado de la fibra se escribirá, en el orden de sus campos pertinentes,
\[
 x=(a;q;\tau,o,h;b_{\rm ref};
       R^\Sigma,q^\Sigma,R^\Pi,q^\Pi;c;
       \operatorname{Sig};C,\partial C;p_{\APP},p_{\TRIT};
       \lambda;\operatorname{gr}_{\mathsf M};\operatorname{mem};
       s;\operatorname{Led}).
\]
La notación de actualización \(x[f\leftarrow v]\) cambia sólo el campo
indicado. Para \(h\in\{(\Sigma,\Pi),(\Pi,\Sigma)\}\), defínase
\[
 h_r(h)=
 \begin{cases}
  (\Sigma,\Pi),&r\in\{1,4,7\},\\
  (\Pi,\Sigma),&r\in\{2,5,8\},\\
 h,&r\in\{3,6,9\}.
 \end{cases}
\]
La coordenada completa de hoja y aritmética se abrevia por
\[
 \widehat h(x):=
 \bigl(h(x),R^\Sigma(x),q^\Sigma(x),R^\Pi(x),q^\Pi(x)\bigr).
\]
Sea \(\widehat{\mathcal H}^{\rm cal}_{\widetilde q_{n,r}}\) el conjunto de
las quíntuplas
\[
 \bigl(h,R^\Sigma_{a_{n,r}},q^\Sigma_{a_{n,r}},
          R^\Pi_{a_{n,r}},q^\Pi_{a_{n,r}}\bigr),
 \qquad
 h\in\{(\Sigma,\Pi),(\Pi,\Sigma)\},
\]
cuyos cuatro datos aritméticos son las dos divisiones euclídeas de la célula
APP etiquetada \(a_{n,r}\). Esta definición directa no presupone un estado
en un portador todavía no construido. Sean
\(\mathsf{Hist}^{\rm cal}_{\widetilde q}\) el conjunto de caminos finitos
admisibles del grafo de aristas etiquetadas que terminan en
\(\widetilde q\), incluida la identidad, y
\(\mathsf B^{\rm cal}_{\widetilde q}\) el conjunto de sus sucesiones
ordenadas de extremos. Definimos
\[
 \mathsf S^{\rm cal}_{\widetilde q}:=
 \widehat{\mathcal H}^{\rm cal}_{\widetilde q}
 \times\{o_{\rightarrow},o_{\circ},o_{\leftarrow}\}
 \times K_{\rm Carr}\times
 \mathsf{Hist}^{\rm cal}_{\widetilde q}\times
 \mathsf B^{\rm cal}_{\widetilde q}\times\mathsf M .
\]
La firma calibrada no se postula: es la aplicación de reunión tipada
\[
 \operatorname{Sig}^{\rm cal}_{\widetilde q}:
 \widehat{\mathcal H}^{\rm cal}_{\widetilde q}
 \times\{o_{\rightarrow},o_{\circ},o_{\leftarrow}\}
 \times K_{\rm Carr}\times
 \mathsf{Hist}^{\rm cal}_{\widetilde q}\times
 \mathsf B^{\rm cal}_{\widetilde q}\times\mathsf M
 \longrightarrow\mathsf S^{\rm cal}_{\widetilde q},
\quad
 (\widehat h,o,c,\lambda,\partial C,\mu)
 \longmapsto(\widehat h,o,c,\lambda,\partial C,\mu).
\]
Asimismo, \(\pi_{\mathcal F}\) denota la proyección que elimina únicamente
las dos primeras coordenadas \((a;q)\) de la tupla de estado; conserva
\(\tau,o,h,b_{\rm ref}\), los cuatro datos euclídeos, acarreo, firma,
cilindro, frontera, prefijos, ruta, grado de memoria, memoria, supervivencia y
registro.
La orientación correspondiente queda fijada por
\(o(+1)=o_{\rightarrow}\), \(o(0)=o_{\circ}\) y
\(o(-1)=o_{\leftarrow}\).
La fibra APP calibrada sobre \(\widetilde q_{n,r}\) contendrá la copia etiquetada
\(a_{n,r}\) ya definida. El estado inicial queda fijado sin coordenadas libres:
\[
 \widehat h_\star:=
 \bigl((\Sigma,\Pi),R^\Sigma_{1,9},q^\Sigma_{1,9},
       R^\Pi_{1,9},q^\Pi_{1,9}\bigr).
\]
\begin{equation}
\begin{aligned}
 x_\star:=\bigl(&a_{0,1};\widetilde q_{0,1};
 +1,o_{\rightarrow},(\Sigma,\Pi);\bot;
 R^\Sigma_{1,9},q^\Sigma_{1,9},R^\Pi_{1,9},q^\Pi_{1,9};0;\\
 &\operatorname{Sig}^{\rm cal}_{\widetilde q_{0,1}}
       (\widehat h_\star,o_{\rightarrow},0,
        \operatorname{id}_{\widetilde q_{0,1}},\varnothing,
        \jmath_{k_0}(0_{\mathsf M_{k_0}}));
 \varnothing,\varnothing;
 (a_{0,1}),(+1);\operatorname{id}_{\widetilde q_{0,1}};
 k_0;0_{\mathsf M_{k_0}};\top;\varnothing\bigr).
\end{aligned}
\label{hmtfg:eq:estado-inicial-calibrado-ch}
\end{equation}
La tupla no contiene coordenadas indeterminadas y satisface
\(s(x_\star)=\top\). Sea \(G_0:=\{x_\star\}\).
Supuesto construido \(G_n\), fijemos \(x\in G_n\),
\(\varepsilon\in\{-1,0,+1\}\) y
\(x_{\varepsilon,0}(x):=x\). Para \(r=1,\ldots,9\), se construye primero
el registro elemental
\begin{equation}
\begin{aligned}
 \ell_{\varepsilon,n,r}:=\bigl(&e_{\varepsilon,n,r},a_{n,r},a_{(n,r)^+},
 \widetilde q_{n,r},\widetilde q_{n,r+1},M_{\rm ph}(r),\\
 &h_r(h(x_{\varepsilon,r-1}(x))),o(M_{\rm ph}(r)),\mathfrak p_\varepsilon,
 d_{\varepsilon,r},\\
 &\mathsf C_{\mathsf M}(e_{\varepsilon,n,r}),
 \kappa_{\mathsf M}(e_{\varepsilon,n,r}),
 \operatorname{Carr}(e_{\varepsilon,n,r})\bigr),
\end{aligned}
 \label{hmtfg:eq:registro-elemental-ext-ch}
\end{equation}
y se pone
\(\operatorname{Led}^{\rm new}_{\varepsilon,n,r}:=
\operatorname{Led}(x_{\varepsilon,r-1}(x))^\frown
\ell_{\varepsilon,n,r}\). Se construye entonces la tupla
\begin{equation}
 y_{\varepsilon,r}(x):=
 x_{\varepsilon,r-1}(x)
 [\tau\leftarrow M_{\rm ph}(r),\
  h\leftarrow h_r(h(x_{\varepsilon,r-1}(x))),\
  o\leftarrow o(M_{\rm ph}(r)),\
  b_{\rm ref}\leftarrow\mathfrak p_\varepsilon],
 \label{hmtfg:eq:tupla-sel-cal-ch}
\end{equation}
donde la selección deja intactos célula, cocientes, acarreo, ruta, memoria,
cilindro, frontera, supervivencia y registro; la coordenada
\(b_{\rm ref}\) hace explícita la reducción balanceada escogida. No se presupone aquí
pertenencia a una relación ambiental.

A continuación se define \(z_{\varepsilon,r}(x)\) campo por campo: su
configuración es \(\widetilde q_{n,r+1}\); su célula completa es
\(a_{(n,r)^+}\), con los
cuatro residuos y cocientes recalculados por las dos divisiones euclídeas;
conserva la hoja seleccionada y la orientación; y efectúa las seis
actualizaciones siguientes, en las que sólo se omite el argumento común \(x\):
\begin{equation}
\begin{aligned}
 C(z_{\varepsilon,r})&:=C(y_{\varepsilon,r})^\frown
       e_{\varepsilon,n,r},\\
 \partial C(z_{\varepsilon,r})&:=\partial C(y_{\varepsilon,r})^\frown
       (\widetilde q_{n,r},\widetilde q_{n,r+1}),\\
 p_{\APP}(z_{\varepsilon,r})&:=p_{\APP}(y_{\varepsilon,r})^\frown
       a_{(n,r)^+},\\
 p_{\TRIT}(z_{\varepsilon,r})&:=p_{\TRIT}(y_{\varepsilon,r})^\frown
       \tau(r),\\
 \operatorname{Led}(z_{\varepsilon,r})&:=
       \operatorname{Led}^{\rm new}_{\varepsilon,n,r},\\
 \lambda(z_{\varepsilon,r})&:=e_{\varepsilon,n,r}
       \lambda(y_{\varepsilon,r}),
\end{aligned}
 \label{hmtfg:eq:tupla-tra-cal-ch}
\end{equation}
y mantiene todavía los campos
\(c,\operatorname{gr}_{\mathsf M},\operatorname{mem},s\) de
\(y_{\varepsilon,r}(x)\). Su firma, en cambio, se vuelve a tipar en la
configuración de destino mediante la prefirma transportada
\begin{equation}
\begin{aligned}
 \operatorname{Sig}(z_{\varepsilon,r}(x))
  :=\operatorname{Sig}^{\rm cal}_{\widetilde q_{n,r+1}}
       \Bigl(&\widehat h(z_{\varepsilon,r}(x)),
             o(z_{\varepsilon,r}(x)),c(y_{\varepsilon,r}(x)),
             \lambda(z_{\varepsilon,r}(x)),\\[-2pt]
             &\partial C(z_{\varepsilon,r}(x)),
             \jmath_{d_{n,r}}
                (\operatorname{mem}_{d_{n,r}}
                   (y_{\varepsilon,r}(x)))\Bigr).
\end{aligned}
\label{hmtfg:eq:prefirma-transporte-cal-ch}
\end{equation}
Así, y sólo en
\(\mathsf{Tra}^{\rm cal}\), se añaden la arista, sus extremos, la marca
\((\mathfrak p_\varepsilon,r,d_{\varepsilon,r})\) y sus cociclos al registro de
incidencias.

El estado inicial satisface
\[
 c(x_\star)=0=\operatorname{Carr}
   (\operatorname{id}_{\widetilde q_{0,1}})
   =\operatorname{Carr}(\lambda(x_\star)).
\]
En los estados completos \(x_{\varepsilon,r-1}(x)\), y también después de la
selección \(y_{\varepsilon,r}(x)\), se conserva el invariante tipado
\(c(v)=\operatorname{Carr}(\lambda(v))\), pues
\(\mathsf{Sel}^{\rm cal}\) no modifica ninguno de esos dos campos. El
transporte avanza la ruta pero deja pendiente el acarreo; por ello la tupla de
preactualización satisface exactamente
\[
 c(z_{\varepsilon,r}(x))=c(y_{\varepsilon,r}(x))
 =\operatorname{Carr}(\lambda(y_{\varepsilon,r}(x))).
\]
La aplicación \(\operatorname{Upd}^{\rm cal}\) restaura entonces
\(c(x_{\varepsilon,r}(x))=
\operatorname{Carr}(\lambda(x_{\varepsilon,r}(x)))\) mediante la ley de
composición, sin atribuir ese invariante al estado intermedio \(z\).

Finalmente se aplica la función unaria
\(\operatorname{Upd}^{\rm cal}_{e_{\varepsilon,n,r}}
=\mathsf{Mem}^{\rm cal}_{e_{\varepsilon,n,r}}\) a esa tupla. Su
resultado \(x_{\varepsilon,r}(x)\) viene dado por
\begin{equation}
\begin{aligned}
 c\bigl(x_{\varepsilon,r}(x)\bigr)
   &=\operatorname{Carr}\bigl(\lambda(z_{\varepsilon,r}(x))\bigr)\\
   &=H_{\rm Carr}(\lambda(y_{\varepsilon,r}(x)))
       \bigl(\operatorname{Carr}(e_{\varepsilon,n,r})\bigr)
     +Y_{\rm Carr}(e_{\varepsilon,n,r})
       \bigl(c(y_{\varepsilon,r}(x))\bigr),\\
 \lambda\bigl(x_{\varepsilon,r}(x)\bigr)
   &=\lambda(z_{\varepsilon,r}(x)),\\
 \operatorname{gr}_{\mathsf M}
      \bigl(x_{\varepsilon,r}(x)\bigr)
   &=d_{n,r}+1,\\
 \operatorname{mem}_{d_{n,r}+1}
       \bigl(x_{\varepsilon,r}(x)\bigr)
   &=\mathsf C_{\mathsf M}(e_{\varepsilon,n,r})
       \operatorname{mem}_{d_{n,r}}(z_{\varepsilon,r}(x))
     +\kappa_{\mathsf M}(e_{\varepsilon,n,r}),\\
 s\bigl(x_{\varepsilon,r}(x)\bigr)
   &=s(e_{\varepsilon,n,r})\wedge s(z_{\varepsilon,r}(x)),\\
 \operatorname{Sig}\bigl(x_{\varepsilon,r}(x)\bigr)
   &=\operatorname{Sig}^{\rm cal}_{\widetilde q_{n,r+1}}
       \Bigl(\widehat h(z_{\varepsilon,r}(x)),
             o(z_{\varepsilon,r}(x)),
             c(x_{\varepsilon,r}(x)),
             \lambda(z_{\varepsilon,r}(x)),
             \partial C(z_{\varepsilon,r}(x)),\\[-2pt]
   &\hspace{112pt}
             \jmath_{d_{n,r}+1}
              \bigl(\operatorname{mem}_{d_{n,r}+1}
                 (x_{\varepsilon,r}(x))\bigr)\Bigr).
\end{aligned}
\label{hmtfg:eq:upd-unaria-tres-vueltas-ch}
\end{equation}
Todos los campos no enumerados en
\eqref{hmtfg:eq:upd-unaria-tres-vueltas-ch} coinciden con los de
\(z_{\varepsilon,r}(x)\). En particular, la firma se calcula con la hoja,
orientación, frontera y ruta ya fijadas por \(\mathsf{Sel}^{\rm cal}\) y
\(\mathsf{Tra}^{\rm cal}\), y con el acarreo y la memoria calculados en las
líneas anteriores; no se define circularmente a partir del estado final.
Ni \(\varepsilon\), ni \(d_{\varepsilon,r}\), ni \(u\) son argumentos
adicionales de \(\operatorname{Upd}^{\rm cal}\): la función lee los tres datos de la
arista ya construida.

\paragraph{Inducción simultánea de los niveles.}
Completadas las nueve etapas anteriores para un \(G_n\) ya construido,
se define inmediatamente
\begin{equation}
 \gamma_{\varepsilon,k_n}(x):=x_{\varepsilon,9}(x),\qquad
 G_{n+1}:=\{\gamma_{\varepsilon,k_n}(x):
       x\in G_n,\ \varepsilon\in\TRIT\}.
 \label{hmtfg:eq:injerto-total-tres-vueltas-ch}
\end{equation}
\label{hmtfg:eq:tres-vueltas-tpk-ch}
La base \(G_0=\{x_\star\}\), las tuplas intermedias y
\eqref{hmtfg:eq:injerto-total-tres-vueltas-ch} constituyen una única definición por
recursión sobre \(n\). Por tanto, ningún portador posterior se utiliza para
suponer la existencia de \(G_n\): el nivel siguiente sólo se introduce
después de construir todas sus aristas desde el nivel precedente.

\paragraph{Relación de extensión TPK y pertenencia efectiva.}
Sea \(\mathcal E^{\rm cal}_{\widetilde q}\) la menor familia etiquetada que contiene
\(x_\star\) y, en cada paso de la recurrencia anterior, las cuatro tuplas
\(x_{\varepsilon,r-1}(x),y_{\varepsilon,r}(x),
z_{\varepsilon,r}(x),x_{\varepsilon,r}(x)\) en sus configuraciones de origen
y destino. Sea asimismo
\(\mathcal A^{\rm cal}_{\widetilde q_{n,r}}:=\{a_{n,r}\}\), y póngase
\[
 \widetilde{\mathcal Q}_{\rm obs}:=
   \{\widetilde q_{n,r}:n\ge0,\ 1\le r\le9\},
 \qquad
 \mathcal E^{\rm cal}:=\bigsqcup_{\widetilde q}
     \mathcal E^{\rm cal}_{\widetilde q},
 \qquad
 \mathcal F^{\rm cal}_{\widetilde q}:=
     \pi_{\mathcal F}(\mathcal E^{\rm cal}_{\widetilde q}).
\]
Para cada profundidad \(d\), sea
\(\mathcal E^{\rm cal}_d:=
\operatorname{gr}_{\mathsf M}^{-1}(d)\subset\mathcal E^{\rm cal}\). La
coordenada explícita de grado, y no la mera pertenencia de un representante a
alguna imagen del sistema directo, hace disjuntos estos portadores. Sobre ellos
quedan tipadas las proyecciones
\[
\begin{aligned}
 \operatorname{cfg}&:\mathcal E^{\rm cal}\longrightarrow
     \widetilde{\mathcal Q}_{\rm obs},\\
 g(x)&:=\operatorname{fase}\!
       \bigl(\pi_{\rm obs}(\operatorname{cfg}(x))\bigr)\in C_9,\\
 c&:\mathcal E^{\rm cal}\longrightarrow K_{\rm Carr},\\
 s&:\mathcal E^{\rm cal}\longrightarrow\{\bot,\top\},\\
 \lambda&:\mathcal E^{\rm cal}_{\widetilde q}
       \longrightarrow\mathsf{Hist}^{\rm cal}_{\widetilde q},\\
 \operatorname{gr}_{\mathsf M}&:\mathcal E^{\rm cal}
       \longrightarrow\mathbb N_{\ge1},\\
 \operatorname{mem}_d&:\mathcal E^{\rm cal}_d
       \longrightarrow\mathsf M_d,\\
 \operatorname{mem}(x)&:=
       \jmath_d\bigl(\operatorname{mem}_d(x)\bigr)\in\mathsf M
       \quad(x\text{ de profundidad }d).
\end{aligned}
\]
Aquí \(\operatorname{cfg}(x)\) es la segunda coordenada de la tupla y
\(\operatorname{fase}\) lee la fase del calendario observable. Para las tres
clases de portador se fijan las diagonales
\[
\begin{aligned}
 \Delta_{\APP,\widetilde q}
   &:=\{(a,a):a\in\mathcal A^{\rm cal}_{\widetilde q}\},\\
 \Delta_{{\rm fib},\widetilde q}
   &:=\{(f,f):f\in\mathcal F^{\rm cal}_{\widetilde q}\},\\
 \Delta_{{\rm enr},\widetilde q}
   &:=\{(x,x):x\in\mathcal E^{\rm cal}_{\widetilde q}\}.
\end{aligned}
\]
En particular,
\[
 a_{0,1}\in\mathcal A^{\rm cal}_{\widetilde q_{0,1}},
 \qquad
 x_\star\in\mathcal E^{\rm cal}_{\widetilde q_{0,1}}.
\]
Sobre estos portadores se
definen por sus grafos completos, y no mediante implicaciones necesarias, las
relaciones
\begin{align}
 \mathsf{Sel}^{\rm cal}_{e_{\varepsilon,n,r}}
   &:=\{(x_{\varepsilon,r-1}(x),y_{\varepsilon,r}(x)):
          x\in G_n\},\\
 \mathsf{Tra}^{\rm cal}_{e_{\varepsilon,n,r}}
   &:=\{(y_{\varepsilon,r}(x),z_{\varepsilon,r}(x)):
          x\in G_n\},\\
 \operatorname{Ext}^{\APP,\rm cal}_{e_{\varepsilon,n,r}}
   &:=\{(a_{n,r},a_{(n,r)^+})\},
 \label{hmtfg:eq:ext-app-calibrada-ch}\\
 \operatorname{Ext}^{\rm fib,cal}_{e_{\varepsilon,n,r}}
   &:=\{(\pi_{\mathcal F}x_{\varepsilon,r-1}(x),
          \pi_{\mathcal F}x_{\varepsilon,r}(x)):x\in G_n\}.
 \label{hmtfg:eq:ext-fibra-calibrada-ch}
\end{align}
La extensión enriquecida sincronizada es
\begin{equation}
 \operatorname{Ext}^{\rm enr,cal}_{e_{\varepsilon,n,r}}
 :=\{(x_{\varepsilon,r-1}(x),x_{\varepsilon,r}(x)):
       x\in G_n\}.
 \label{hmtfg:eq:ext-enriquecida-calibrada-ch}
\end{equation}
La tercera aplicación queda definida por
\begin{equation}
 \operatorname{graph}
 \bigl(\operatorname{Upd}^{\rm cal}_{e_{\varepsilon,n,r}}\bigr)
 :=\{(z_{\varepsilon,r}(x),x_{\varepsilon,r}(x)):
       x\in G_n\},
 \label{hmtfg:eq:grafo-upd-calibrado-ch}
\end{equation}
donde el segundo componente es exactamente el de
\eqref{hmtfg:eq:upd-unaria-tres-vueltas-ch}. Por consiguiente,
\begin{equation}
 \mathcal U^{(1),\rm cal}_{e_{\varepsilon,n,r}}
 :=\operatorname{Upd}^{\rm cal}_{e_{\varepsilon,n,r}}\circ
   \mathsf{Tra}^{\rm cal}_{e_{\varepsilon,n,r}}\circ
   \mathsf{Sel}^{\rm cal}_{e_{\varepsilon,n,r}}
 =\{(x_{\varepsilon,r-1}(x),x_{\varepsilon,r}(x)):
      x\in G_n\}.
 \label{hmtfg:eq:transicion-cal-en-tpk-ch}
\end{equation}
Cada operador está indexado por la arista
\(e_{\varepsilon,n,r}\) producida por la reducción
\(\mathfrak p_\varepsilon\), y no por un trit suministrado desde fuera. La
coordenada \(b_{\rm ref}\) distingue además sus codominios. Por ello
\(\mathsf{Sel}^{\rm cal}_{e}\),
\(\mathsf{Tra}^{\rm cal}\) y \(\operatorname{Upd}^{\rm cal}\) son funciones
sobre sus respectivos dominios, mientras
la unión
\(\bigcup_{\varepsilon\in\TRIT}
\mathcal U^{(1),\rm cal}_{e_{\varepsilon,n,r}}\) conserva deliberadamente
las tres salidas de la trisección.

Para \(\bullet\in\{\APP,{\rm fib},{\rm enr}\}\), las identidades y las
composiciones no quedan implícitas. Se definen
\begin{align}
 \operatorname{Ext}^{\bullet,\rm cal}_{\operatorname{id}_{\widetilde q}}
   &:=\Delta_{\bullet,\widetilde q},\\
 \operatorname{Ext}^{\bullet,\rm cal}_{e_t\cdots e_1}
   &:=\operatorname{Ext}^{\bullet,\rm cal}_{e_t}\circ\cdots\circ
      \operatorname{Ext}^{\bullet,\rm cal}_{e_1}
 \label{hmtfg:eq:ext-palabra-calibrada-ch}
\end{align}
para toda palabra admisible de aristas. Las dos clausuras de caminos APP y
fibra se forman por separado,
\begin{equation}
 \operatorname{Path}(\operatorname{Ext}^{\APP,\rm cal}),\qquad
 \operatorname{Path}(\operatorname{Ext}^{\rm fib,cal}),
 \label{hmtfg:eq:dos-clausuras-ext-ch}
\end{equation}
y \(\mathsf{Tra}^{\rm cal}\) las sincroniza arista por arista.

Finalmente, los niveles de estados completos y sus truncamientos se fijan
por la propia recurrencia:
\begin{align}
 H^{\rm cal}_{d_{n,r}}
   &:=\{x_{\varepsilon,r-1}(x):
          x\in G_n,\ \varepsilon\in\TRIT\},\\
 H^{\rm cal}_{d_{n,r}+1}
   &:=\{x_{\varepsilon,r}(x):
          x\in G_n,\ \varepsilon\in\TRIT\},\\
 \rho^{\rm cal}_{d_{n,r}+1,d_{n,r}}
   \bigl(x_{\varepsilon,r}(x)\bigr)
   &:=x_{\varepsilon,r-1}(x).
 \label{hmtfg:eq:truncamiento-calibrado-definido-ch}
\end{align}
El último registro conserva \((n,r,\mathfrak p_\varepsilon)\), de modo que
el predecesor del miembro izquierdo es único y \(\rho^{\rm cal}\) está bien
definida. Para \(b>a\), se pone
\(\rho^{\rm cal}_{b,a}:=\rho^{\rm cal}_{a+1,a}\circ\cdots\circ
\rho^{\rm cal}_{b,b-1}\). Esta aplicación recupera el estado anterior
completo, incluida su firma; no pretende invertir una coordenada sobrescrita
sin consultar el registro.

\begin{proposition}[Realización autónoma del TPK mediante transiciones explícitas]
\label{hmtfg:prop:modelo-tpk-calibrado-ch}
La familia
\[
\begin{aligned}
 \mathfrak T_{\rm TPK}^{\rm cal}:=\bigl(
 &(\widetilde q_{n,r})_{n,r},\pi_{\rm obs},
   \widetilde F_{\rm obs},M_{\rm ph},\\
 &(\mathcal A_{\widetilde q}^{\rm cal})_{\widetilde q},
   (\mathcal E_{\widetilde q}^{\rm cal})_{\widetilde q},
   (\mathcal F_{\widetilde q}^{\rm cal})_{\widetilde q},
   \operatorname{Ext}^{\APP,\rm cal},
   \operatorname{Ext}^{\rm fib,cal},
   \operatorname{Ext}^{\rm enr,cal},\rho^{\rm cal},\\
 &\mathsf{Sel}^{\rm cal},\mathsf{Tra}^{\rm cal},
   \operatorname{Upd}^{\rm cal},\mathsf C_{\mathsf M},
   \kappa_{\mathsf M},\operatorname{gr}_{\mathsf M},
   H_{\rm Carr},Y_{\rm Carr},\\
 &\operatorname{Sig}^{\rm cal},\pi_{\mathcal F}\bigr).
\end{aligned}
\]
es un modelo enriquecido del TPK. En particular, cada par del miembro derecho
de \eqref{hmtfg:eq:transicion-cal-en-tpk-ch} es una arista TPK efectiva, y no la
consecuencia inferida de una condición solamente necesaria.
\end{proposition}

\begin{proof}
Las etiquetas \((n,r)\) hacen disjuntas las copias de los portadores y fijan
sin ambigüedad cada origen y destino. La relación
\(\operatorname{Ext}^{\APP,\rm cal}\) transporta la célula completa y
preserva las dos identidades euclídeas; la relación
\(\operatorname{Ext}^{\rm fib,cal}\) relaciona las fibras de los estados
completos anterior y posterior: transporta hoja, orientación, prefijo,
cilindro, frontera y registro e incorpora el acarreo, la memoria y la firma
actualizados. La selección no modifica el registro y el transporte concatena
exactamente una entrada y produce la prefirma de destino
\eqref{hmtfg:eq:prefirma-transporte-cal-ch}. La actualización recalcula la firma mediante la
aplicación tipada \(\operatorname{Sig}^{\rm cal}_{\widetilde q}\) desde
hoja, orientación, acarreo, ruta, frontera y memoria ya determinados. La aplicación
\(\operatorname{Upd}^{\rm cal}\) es unaria porque la arista contenida en
\(z_{\varepsilon,r}(x)\) determina sus cociclos.

Las identidades son las diagonales y las extensiones por palabras son las
composiciones de \eqref{hmtfg:eq:ext-palabra-calibrada-ch}. La concatenación de
prefijos y registros es asociativa; la memoria satisface la ley de cociclo y
el acarreo satisface
\eqref{hmtfg:eq:ley-carry-correcta-ch}. Por tanto, las dos parentizaciones de tres
aristas producen el mismo estado. La fórmula
\eqref{hmtfg:eq:truncamiento-calibrado-definido-ch} define el truncamiento sobre
los estados completos y la unicidad del último registro prueba que recupera
el origen en todos los campos. Éstas son las leyes de identidad, fuente,
destino, composición, memoria, acarreo y truncamiento del TPK enriquecido.
En particular, las relaciones de extensión quedan determinadas como grafos
de aplicaciones efectivas y no como consecuencias de meras condiciones
necesarias.
\end{proof}

Cada \(\gamma_{\varepsilon,k_n}\) comprende exactamente nueve aplicaciones
de \(\operatorname{Upd}^{\rm cal}_{e_{\varepsilon,n,r}}\circ
\mathsf{Tra}^{\rm cal}_{e_{\varepsilon,n,r}}\circ
\mathsf{Sel}^{\rm cal}_{e_{\varepsilon,n,r}}\). La
supervivencia se probará mediante una cola construida y no se utilizará para
definir \(G_n\).

\begin{lemma}[Trisección interna efectiva de una vuelta]
\label{hmtfg:lem:elevacion-catalogal-nonadica}
Para todo \(n\ge0\), todo \(x\in G_n\) y todo
\(\varepsilon\in\{-1,0,+1\}\), la expresión
\(\gamma_{\varepsilon,k_n}(x)\) está definida por nueve transiciones TPK,
con sus relaciones sincronizadas
\(\operatorname{Ext}^{\APP,\rm cal}\) y
\(\operatorname{Ext}^{\rm fib,cal}\), y satisface
\begin{align}
 \rho^{\rm cal}_{k_n+9,k_n}\gamma_{\varepsilon,k_n}(x)&=x,
 \label{hmtfg:eq:truncamiento-injerto-ch}\\
 g\bigl(\gamma_{\varepsilon,k_n}(x)\bigr)
   &=g(x),&
 \jmath_{k_n+9}\!\left(\operatorname{mem}_{k_n+9}
        (\gamma_{\varepsilon,k_n}(x))\right)
   &=\jmath_{k_n}\!\left(\operatorname{mem}_{k_n}(x)\right)+u,
 \label{hmtfg:eq:fase-memoria-injerto-ch}\\
 c\bigl(\gamma_{\varepsilon,k_n}(x)\bigr)
   &=c(x)+\varepsilon\kappa_{\rm Carr}.&&
 \label{hmtfg:eq:carry-injerto-ch}
\end{align}
\label{hmtfg:eq:tres-vueltas-propiedades-ch}
\label{hmtfg:eq:tres-vueltas-acarreo-ch}
El registro ordenado de las nueve aristas determina \(\varepsilon\) de manera única.
Las tres imágenes son, por tanto, disjuntas. Cada extremo pertenece a
\(G_{n+1}\) y vuelve a admitir las tres aplicaciones
\(\gamma_{-1,k_{n+1}},\gamma_{0,k_{n+1}},
\gamma_{+1,k_{n+1}}\).
\end{lemma}

\begin{proof}
Las fórmulas \eqref{hmtfg:eq:celula-completa-fase-ch} recorren una vez las nueve
fases y devuelven la copia \(a_{n+1,1}\) de \(a(1)\). En cada paso, la relación APP transporta el
séxtuplo completo; \(\mathsf{Sel}^{\rm cal}\) fija el tipo TRIT desde la marca de fase
\(r\) conservada en el estado;
y \(\mathsf{Tra}^{\rm cal}\) añade exactamente un símbolo al prefijo, un tramo al
cilindro y su par ordenado de extremos a la frontera. Por inducción sobre
\(r\), las tuplas explícitas \(y_{\varepsilon,r}(x)\) y
\(z_{\varepsilon,r}(x)\) satisfacen los predicados de las relaciones
indicadas, y la actualización unaria produce el estado del nivel siguiente.
La identidad
\(\rho^{\rm cal}_{d_{n,r}+1,d_{n,r}}
(x_{\varepsilon,r}(x))=x_{\varepsilon,r-1}(x)\) es la definición
\eqref{hmtfg:eq:truncamiento-calibrado-definido-ch}; su buena definición procede de
la unicidad del último registro. La composición de las nueve aplicaciones de
truncamiento recupera \(x\), lo que prueba
\eqref{hmtfg:eq:truncamiento-injerto-ch} sin invertir por separado los campos
sobrescritos.

La fase \(9\to1\) aparece una sola vez. Por
\eqref{hmtfg:eq:memoria-arista-ch} y la ley de cociclo, la parte lineal de memoria
transporta la memoria previa y su parte traslacional suma exactamente \(u\)
en el límite directo. La compatibilidad
\(\jmath_{d+1}\jmath_{d+1,d}=\jmath_d\) prueba
\eqref{hmtfg:eq:fase-memoria-injerto-ch}; el retorno concierne a la fase, no al
estado enriquecido.

Por \eqref{hmtfg:eq:carry-arista-derivado-ch}, ocho aristas poseen incremento nulo y
la arista que completa \(\mathfrak p_\varepsilon\) posee el incremento calculado
\(\chi(\varepsilon,\varepsilon)\kappa_{\rm Carr}\). La ley
\eqref{hmtfg:eq:ley-carry-correcta-ch} y
\eqref{hmtfg:eq:carry-no-parametrico-ch} dan
\eqref{hmtfg:eq:carry-injerto-ch}. El registro conserva
\((\mathfrak p_\varepsilon,2\varepsilon,r_3(2\varepsilon),
\chi(\varepsilon,\varepsilon))\); los tres valores de este registro son
distintos incluso cuando el incremento escalar es cero. Por ello las imágenes
no se identifican y el decodificador de bloque recupera una única
\(\varepsilon\).

Las definiciones de las aristas dependen sólo de la nueva fase y concatenan,
en vez de sustituir, prefijo, cilindro, frontera y registro. También son
invariantes bajo
\(\jmath_d\operatorname{mem}_d\mapsto
  \jmath_d\operatorname{mem}_d+nu\)
y bajo la traslación del acarreo de
entrada. En consecuencia, la misma construcción se aplica al extremo de
cualquier vuelta. Para cada estado finito, la sucesión explícita
\[
 x,\quad\gamma_{0,k_n}(x),\quad
 \gamma_{0,k_{n+1}}\gamma_{0,k_n}(x),\quad\ldots
\]
es una cola infinita compatible. Esto demuestra la supervivencia de todos los
estados construidos y la disponibilidad reiterada de las tres ramas, sin
postular ninguna de ambas propiedades. Por consiguiente podemos, finalmente,
identificar, de acuerdo con las definiciones de nivel anteriores,
\begin{equation}
 H_{k_n}^{\rm cal}:=G_n
 \label{hmtfg:eq:subarbol-calibrado-superviviente-ch}
\end{equation}
para todo \(n\ge0\).
\end{proof}

El cierre anterior distingue los tres relojes implicados. Nueve emisiones
devuelven la fase en \(C_9\) y avanzan la memoria. Doce vueltas, es decir,
\(108\) emisiones, devuelven además la posición del calendario observable;
tampoco entonces se borran el registro ni la memoria. El operador
\(\mathcal K_{\rm ph}\) actúa únicamente después, como reconocimiento de
memoria reducida de las vueltas ya construidas; nunca genera las aristas de
\(\operatorname{Ext}^{\rm enr,cal}\).
